\section{Method}\label{sec:method}

\subsection{Preliminaries}\label{sec:background}

\paragraph{Goal Accessibility Ratio.}
Standard metrics such as perplexity and attention entropy capture generation quality but not \emph{what} the model attends to. In multi-turn dialogue, a model may generate fluent text while progressively ignoring task specifications in the system prompt. We need a metric that directly measures whether goal-relevant context receives attention. To formalize this, we denote the task specification tokens in the system prompt as \emph{goal tokens} $\mathcal{G}$. \citet{dongre2026attention} defined the Goal Accessibility Ratio (GAR) for a model with $L$ layers and $H$ heads at turn $\tau$ as
\begin{equation}\label{eq:gar}
  \text{GAR}(\tau) = \frac{1}{L \cdot H \cdot |\mathcal{R}_\tau|} \sum_{\ell=1}^{L} \sum_{h=1}^{H} \sum_{i \in \mathcal{R}_\tau} \sum_{j \in \mathcal{G}} A_{i,j}^{(\ell,h)},
\end{equation}
where $\mathcal{R}_\tau$ is the set of response tokens at turn $\tau$ and $A_{i,j}^{(\ell,h)}$ is the attention weight from token $i$ to token $j$ at layer $\ell$, head $h$.

\paragraph{Two GAR variants.}
\textbf{GAR-Dongre} is the original definition (Eq.~\ref{eq:gar}), averaging over all response tokens. However, averaging dilutes the signal with intermediate positions that may legitimately attend to local context rather than the system prompt. The last token of each assistant turn is where the model aggregates information for the next continuation; attention at this position most directly reflects the model's awareness of task specifications. \textbf{GAR-ours} restricts the measurement accordingly:
\begin{equation}\label{eq:gar-ours}
  \text{GAR\text{-}ours}(\tau) = \frac{1}{L \cdot H} \sum_{\ell=1}^{L} \sum_{h=1}^{H} \sum_{j \in \mathcal{G}} A_{t_\tau,\, j}^{(\ell,h)},
\end{equation}
where $t_\tau$ is the last token index at turn $\tau$. Attention at this position reflects the information drawn upon for continuation; we use GAR-ours as the optimization target for $\mathcal{L}_\text{GAR}$ (\S\ref{sec:loss}). Empirically, GAR-Dongre is diluted by NTP attention patterns to ${\sim}$0.01--0.02 in multi-turn settings, losing discriminative power across configurations; GAR-ours retains an order-of-magnitude range (0.019--0.50). The contrast itself is informative: goal attention concentrates at the last-token position, and only this concentrated pattern is vulnerable to distribution shift (\S\ref{sec:analysis}). Throughout this work, GAR serves primarily as a \emph{diagnostic probe}: it measures whether auxiliary attention training transfers to autoregressive generation, rather than being an end-to-end optimization objective whose increase directly improves downstream task performance (the correlation with instruction retention is weak; $r$\,=\,0.200, $p$\,=\,0.007; \S\ref{sec:discussion}).

\subsection{ThreadGate Architecture}

\begin{figure}[t]
  \centering
  \includegraphics[width=\columnwidth]{fig4_attention_heatmap.pdf}
  \caption{Per-turn GAR-ours under AR generation ($N$\,=\,5 conversations, 10 turns). All fine-tuned variants cluster together (77--81\% decay; $\Delta$\,=\,3.3\,pp between ThreadGate and LoRA-only); Base decays 87.5\%. Gating adds zero incremental effect.}
  \label{fig:attention-heatmap}
\end{figure}

Different transformer layers serve different roles: early layers capture local syntax while later layers integrate long-range context \citep{gromov2024unreasonable}. A global intervention would uniformly scale all layers, but goal attention decay is layer-dependent: some layers contribute more to cross-turn information routing than others. ThreadGate therefore adds a \emph{per-layer} gating mechanism to the residual stream, allowing the model to learn which layers to modulate. At each layer $l$, the residual update is scaled by a scalar gate:
\begin{equation}
  x_{l+1} = x_l + g_l \cdot f_l(x_l),
\end{equation}
where $f_l$ denotes the sublayer computation (self-attention or feed-forward network, FFN). Since residual networks can be viewed as ensembles of shallow paths \citep{veit2016residual}, per-layer gating selectively modulates individual paths within this ensemble. This resembles ReZero \citep{bachlechner2021rezero}, which initializes a single trainable scalar at zero to stabilize deep network training, but ThreadGate differs in three ways: (i)~gates are \emph{content-conditioned}, a function of the hidden state rather than a fixed scalar; (ii)~gates are \emph{turn-conditioned}, incorporating conversation-position features; (iii)~gates are initialized near identity ($\sigma(1.5) \approx 0.82$) rather than zero, preserving the pretrained model's computation.

Concretely, each gate $g_l$ is produced by a two-layer MLP that takes both content and turn features as input:
\begin{align}\label{eq:gate}
  g_l = \sigma\bigl(W_2\,\text{ReLU}(&W_1\,[\text{proj}(x_l);\;\text{emb}(t_d, t_i)] \notag\\
  &+ b_1) + b_2\bigr),
\end{align}
where $\text{proj}(\cdot)$ reduces the hidden dimension to 16, and $\text{emb}(t_d, t_i)$ provides a 16-dimensional turn embedding encoding turn distance $t_d$ (turns since system prompt) and turn index $t_i$ via learned lookup tables. The 32-dimensional input passes through a 64-unit hidden layer to a scalar gate in $(0, 1)$ via sigmoid. For Qwen-2.5-7B ($L\!=\!28$), this adds ${\sim}$57K parameters ($<$0.001\% of 7B). To isolate the contribution of turn conditioning, the \textbf{Uniform Gate} ablation replaces $\text{emb}(\cdot)$ with a learnable constant, removing turn awareness while retaining content-conditioned gating. To preserve pretrained behavior, gate biases are initialized to 1.5 (${\sim}$82\% pass-through) rather than zero as in ReZero, and gate parameters receive a $10\times$ learning rate multiplier to accelerate adaptation.

\subsection{Training Objective}\label{sec:loss}

Next-token prediction alone provides no gradient signal for attention allocation; a model can achieve low perplexity while distributing attention arbitrarily. The auxiliary GAR loss provides explicit pressure to maintain goal-token attention, while a diversity regularizer prevents gate collapse. The training objective combines three terms:
\begin{equation}
  \mathcal{L} = \mathcal{L}_\text{NTP} + \lambda_\text{GAR}\,\mathcal{L}_\text{GAR} + \lambda_\text{div}\,\mathcal{L}_\text{div}.
\end{equation}
$\mathcal{L}_\text{NTP}$ is the standard next-token prediction loss. For $\mathcal{L}_\text{GAR}$, we use MSE rather than cross-entropy because GAR is a continuous ratio and MSE provides \emph{bidirectional} pressure, penalizing both under- and over-attention and preventing identity collapse (Appendix~\ref{sec:appendix-loss}):
\begin{equation}
  \mathcal{L}_\text{GAR} = \text{MSE}(\text{GAR\text{-}ours}, \tau),
\end{equation}
where $\tau = 0.5$ targets allocating half of the last token's attention to goal tokens. Without diversity pressure, all gates converge near 1.0; the gate diversity loss prevents this:
\begin{equation}
  \mathcal{L}_\text{div} = \text{ReLU}(\sigma_\text{min} - \text{std}(\mathbf{g})),
\end{equation}
where $\mathbf{g} = [g_1, \ldots, g_L]$ and $\sigma_\text{min} = 0.05$ is a minimum standard deviation threshold ($\lambda_\text{GAR} = \lambda_\text{div} = 1.0$). This formulation emerged from seven iterations spanning four loss families (Appendix~\ref{sec:appendix-loss}): cumulative (v1--v2) and maximize (v3--v6) objectives either starve gate gradients at equilibrium or drive gates toward identity; only MSE provides the bidirectional pressure needed to prevent identity collapse.

\subsection{Controlled Comparison Design}\label{sec:comparison-design}

We design a five-way comparison where each adjacent pair isolates one factor: \textbf{ThreadGate (TG) vs.\ Uniform Gate (UG)} (turn conditioning), \textbf{TG vs.\ LoRA-only} (gating mechanism), \textbf{ThreadGate vs.\ +PCGrad} (gradient conflict via gradient surgery \citep{yu2020pcgrad}), and \textbf{Fine-tuned vs.\ Base} (LoRA adaptation). All fine-tuned variants share training data, LoRA configuration, optimizer, and seed (Appendix~\ref{sec:appendix-hyper}).

